% 合成示例（synthetic example），非正式案例结果
% Repaired version of bad_project (repair decisions made by the Agent;
% scripts never auto-edit):
%   1) vaswani2017attentoin -> vaswani2017attention (strong candidate, typo)
%   2) zhang2020residual has no candidate and no real source: the sentence
%      actually concerns residual networks, so it now cites the verified
%      he2016deep entry; the phantom key disappears
%   3) second \label{sec:intro} -> \label{sec:background}
%   4) he2016deep gained year = {2016} (see refs.bib)
%   5) knuth1984texbook was a missing citation: cited once in the text
%      (entry kept, not deleted)
\documentclass{article}
\usepackage{amsmath}
\title{A Mini Survey Fragment on Sequence Modeling}
\author{Example Author}
\date{}
\begin{document}
\maketitle

\section{Introduction}\label{sec:intro}
The Transformer architecture~\cite{vaswani2017attention} changed sequence
modeling, and the original paper gives full details. For deep learning
fundamentals see \cite{goodfellow2016deep,lecun2015deep}. Another
work~\cite{he2016deep} discusses training stability of very deep networks;
the attention computation is shown in Eq.~\eqref{eq:attn}.

\section{Background}\label{sec:background}
\begin{equation}\label{eq:attn}
  \mathrm{Attention}(Q,K,V) =
  \mathrm{softmax}\!\left(\frac{QK^{T}}{\sqrt{d_k}}\right)V
\end{equation}
Residual networks~\cite{he2016deep} ease the training of deep models.

\section{Conclusion}
See Section~\ref{sec:intro} for more background; the typesetting toolchain
used here is documented in \cite{knuth1984texbook}.

\bibliographystyle{plain}
\bibliography{refs}
\end{document}
